\section{Motivation and practical value}
A frozen LTI model changes when the same vehicle is observed at different speeds.
Clustering those models without accounting for speed can therefore confuse an
operating point with a vehicle identity. An LPV description encodes that these
observations belong to one structural system evaluated at different values of a
known scheduling variable.

The distinction may reduce model and controller proliferation. A purely LTI
architecture can require approximately
\begin{equation}
  N_{\mathrm{families}}N_{\mathrm{speed\ regions}}
\end{equation}
calibrated controllers, whereas the proposed representation requires roughly one
scheduled controller family per relevant structural class. This matters for
calibration, deployment, validation, maintenance, and eventual certification.

Federation becomes particularly useful when individual clients cover only parts
of the scheduling range. Several compatible clients can jointly identify an LPV
family over an envelope that no single client has explored, without transferring
raw trajectories. Structural specialization simultaneously limits negative
transfer from physically incompatible clients. The full architecture separates
three questions:
\begin{center}
\begin{tabular}{lll}
\toprule
Quantity & Meaning & Mechanism\\
\midrule
$\rho(t)$ & current operating condition & LPV scheduling\\
$c_i$ & persistent dynamics family & clustering/personalization\\
$\Sigma_i$ & reliability of local evidence & uncertainty-aware federation\\
\bottomrule
\end{tabular}
\end{center}

This work continues the uncertainty-aware clustered identification and LQI design
pipeline developed for the IFAC study. Its new contribution is to give scheduling
variation and persistent fleet heterogeneity distinct mathematical roles.

